\begin{table}[!htbp]
\centering
\caption{Positioning of the present review relative to recent geological CO$_2$ storage and ML reviews.}
\label{tab:positioning}
\small
\setlength{\tabcolsep}{4pt}
\begin{tabularx}{\textwidth}{>{\RaggedRight\arraybackslash}p{0.16\textwidth} >{\RaggedRight\arraybackslash}p{0.18\textwidth} >{\RaggedRight\arraybackslash}p{0.25\textwidth} >{\RaggedRight\arraybackslash}X}
\toprule
Review & Principal scope & Existing contribution & Distinction of the present review \\
\midrule
Yan et al. (2021) & ML across CCUS & Surveys machine-learning applications across the wider CCUS chain. & Restricts the evidence synthesis to storage decisions and their validation requirements. \\
Li et al. (2023) & Storage mechanisms, projects, and ML & Integrates storage mechanisms, project experience, and emerging ML applications. & Separates primary ML evidence from mechanism, project-context, and benchmark-design anchors. \\
Mao and Jahanbani Ghahfarokhi (2024) & Intelligent decision-making for geological storage & Reviews intelligent decision-making across geological storage workflows. & Uses a verified, missingness-aware corpus to qualify decision claims. \\
Li et al. (2024) & Statistical and computational geological storage & Synthesizes statistical and computational methods, including uncertainty treatment. & Links evidence maturity to seven decision-specific minimum validation packages. \\
Marvin et al. (2025) & Progress and implications of ML in geological storage & Connects ML progress to monitoring, uncertainty, anomaly detection, and application needs. & Traces every quantitative element to verified records while retaining missing evidence explicitly. \\
Lin et al. (2025) & Comprehensive ML-driven geological storage & Provides broad coverage of ML methods and geological-storage applications. & Uses non-compensatory packages to distinguish demonstrations from broader decision support. \\
\bottomrule
\end{tabularx}
\end{table}
